\begin{opinion}
    El presente trabajo de diploma aborda un problema de elevada pertinencia científica y sanitaria para Cuba: la predicción de la eficacia del candidato vacunal Quimi-Vio\textregistered 11, destinado a la protección de la población adulta entre 50 y 74 años frente a la enfermedad neumocócica, a partir de la eficacia ya demostrada de la vacuna pediátrica Quimi-Vio\textregistered 7. La autora, Marian Aguilar Tavier, asumió este reto con rigor metodológico, capacidad de estudio independiente y una sólida integración de conocimientos de inmunología, estadística y programación, cualidades que se reflejan tanto en el desarrollo de la investigación como en el resultado final que hoy se presenta.
    
    Desde el punto de vista metodológico, el trabajo se distingue por la aplicación correcta y bien fundamentada del enfoque de inmunopuente (immunobridging), una metodología de frontera en el desarrollo de vacunas que permite inferir la eficacia de un nuevo candidato vacunal sin necesidad de un ensayo clínico de eficacia a gran escala. La estudiante demuestra dominio tanto de los fundamentos inmunológicos del problema, como de las herramientas estadísticas necesarias para abordarlo con el rigor que exige la investigación regulatoria en vacunas.
    
    Particularmente meritorio resulta el tratamiento dado a la censura de los datos, un fenómeno frecuente y técnicamente complejo en los ensayos de inmunogenicidad, que en este estudio alcanzó porcentajes considerables en varios serotipos. La comparación sistemática entre el método de sustitución simple, ampliamente utilizado, pero con limitaciones teóricas conocidas, y el modelo de regresión Tobit estimado por máxima verosimilitud, constituye un aporte metodológico de valor, pues demuestra de forma cuantitativa cómo la elección del método estadístico puede alterar la magnitud de las estimaciones y la amplitud de los intervalos de confianza, y por tanto la validez misma de las conclusiones regulatorias derivadas del estudio.
    
    El trabajo evidencia además una notable capacidad de análisis crítico al interpretar los resultados: la autora no se limita a reportar que siete de los nueve serotipos comunes cumplieron el criterio de no inferioridad, sino que profundiza en las causas de la excepción observada en los serotipos 14 y 23F, vinculándola correctamente con una mayor respuesta inmunológica de referencia en la población pediátrica y con la variabilidad de los datos. Esta capacidad de ir más allá del resultado numérico para ofrecer una interpretación biológica y estadísticamente coherente es propia de un trabajo de calidad.
    
    Durante el desarrollo de la tesis, la estudiante mostró iniciativa, capacidad crítica para interpretar resultados y disposición constante para incorporar las observaciones realizadas. Demostró además una gran autonomía en la búsqueda bibliográfica y en la asimilación de conceptos avanzados, logrando un balance óptimo entre la teoría estadística y su aplicación práctica. Marian es una joven sumamente dinámica y activa; su excelente capacidad de expresión y su elocuencia natural se complementaron a la perfección con su disciplina y constancia. Haber sido su profesora desde el segundo año de la carrera y guiarla hasta este ejercicio de culminación de estudios ha sido un verdadero privilegio. Su receptividad ante las críticas y su compromiso con la excelencia permitieron una evolución constante del manuscrito, haciendo de este proceso de tutoría una experiencia sumamente gratificante al ver su notable crecimiento académico y personal.
    
    Por todo lo anteriormente expuesto, consideramos que el trabajo de diploma titulado “Predicción de eficacia del candidato vacunal Quimi-Vio\textregistered 11 en adultos” cumple con los requisitos científicos y metodológicos exigidos para esta categoría de trabajo, y recomendamos su presentación ante el tribunal correspondiente para la defensa pública, otorgándole la calificación de Excelente (5 puntos).
    
    \vspace{4cm}
    \begin{flushright}
    	\begin{minipage}{0.5\textwidth}
    		\rule{0.8\linewidth}{0.4pt}\\[4pt]
    		\raggedleft
    		MSc. Yanet García Serrano.\\
    		Facultad de Matemática y Computación\\
    		Universidad de la Habana\\
    		Junio, 2026
    	\end{minipage}
    \end{flushright}
\end{opinion}